\section{引言}

\subsection{应用背景与预测需求}

药物在特定溶剂和温度下的溶解度影响溶剂筛选、结晶工艺和实验条件选择。候选药物、溶剂与温度的组合数远大于能够逐一测量的条件数，因此可以用已知结构和实验条件估计对数溶解度，为安排测量顺序提供依据\cite{vermeire2022,fastsolv2025}。本文关注这一任务中的推广环节：给定训练阶段未见的药物母体，以及一个溶剂和温度，预测其条件溶解度。

\subsection{数据挑战与研究问题}

一条观测对应一个具体的药物--溶剂--温度条件。同一药物母体可贡献多条观测，且不同母体的记录数相差较大。按行划分会使同一母体的记录可能同时出现在训练与验证数据中\cite{leakage2023}，逐行平均也会使高频母体在整体误差中占更大比重。为与实际预测对象对齐，本文在数据划分和主要误差汇总中都采用药物母体作为基本单位。

基于该评价设计，本文依次考察三个问题。RQ1 关注分子表示和训练设计，比较基础 ExtraTrees 与使用丰富描述符、显式交互和候选加权策略的树模型。RQ2 衡量融合增量，检验 TabPFN 能否补充树分支，以及收益在不同母体和配置间如何变化。RQ3 关注模型依据与适用范围，分析高误差条件、分支输出和温度响应，并考察实测数据中的组合内温度关联。

\subsection{本文工作与贡献}

本文将分子描述符和条件特征接入 ExtraTrees 主分支，引入 TabPFN 辅助预测，并以内层母体分组验证选择两条分支的配置与融合权重。随后在共同的外层折上比较完整流程，结合配对差值、分层误差和特征归因说明模型何时受益、何时容易偏离实测值。完整开发集拟合的模型接入交互应用，可输出两条分支的预测、条件响应曲线及开发数据覆盖信息。
